\section{The Cattle of the Sun: Archimedes' challenge}
\label{sec:archimedes-cattle}

The cattle problem attributed to Archimedes asks for a herd belonging to
the Sun, with bulls and cows in four colors: white, black, yellow, and
dappled. Its proportions can be written as linear equations. Two further
conditions ask for a square and a triangle. Together they lead to a Pell
equation with a remarkably large first solution.

This optional application begins with the
\href{https://math.nyu.edu/Archimedes/Cattle/Statement.html}{ancient problem's conditions}.
For a detailed modern exposition, see Ilan Vardi,
\href{https://doi.org/10.1080/00029890.1998.12004887}{``Archimedes' Cattle Problem,''
\emph{American Mathematical Monthly} \textbf{105} (1998), 305--319};
an \href{https://www.maths.tcd.ie/~lebed/Cattle\%20problem\%20Ilan\%20Vardi.pdf}{accessible preprint}
is also available. The calculations below can be reproduced with exact
integer arithmetic in the first notebook.

\subsection*{Eight unknowns, seven linear conditions}
Let $W,B,Y,S$ count white, black, yellow, and spotted (dappled) bulls;
let $w,b,y,s$ count cows of the corresponding colors. The conditions are
\begin{align*}
 W&=\left(\frac12+\frac13\right)B+Y,
 &B&=\left(\frac14+\frac15\right)S+Y,
 &S&=\left(\frac16+\frac17\right)W+Y,\\
 w&=\left(\frac13+\frac14\right)(B+b),
 &b&=\left(\frac14+\frac15\right)(S+s),\\
 y&=\left(\frac16+\frac17\right)(W+w),
 &s&=\left(\frac15+\frac16\right)(Y+y).
\end{align*}
The cow equations use an entire herd of the specified color, including
its bulls. The two additional conditions concern bulls only:
\begin{equation}\label{eq:cattle-shapes}
 W+B=\text{a square},\qquad
 Y+S=\frac{h(h+1)}2\quad\text{for some integer }h\ge1.
\end{equation}

Eliminate $B,Y,S$ from the first three equations to obtain
\[
 B=\frac{267}{371}W,\qquad
 Y=\frac{297}{742}W,\qquad
 S=\frac{790}{1113}W.
\]
The four cow equations form a cycle. Substituting around that cycle
expresses $w$ as a known multiple of $W$ plus
$\frac7{12}\frac9{20}\frac{11}{30}\frac{13}{42}w$.
Since this coefficient is less than $1$, $w$, and then every cow count,
is uniquely determined by $W$. Clearing denominators gives
\begin{center}
\begin{tabular}{lrr}
\toprule
Color & Bulls & Cows\\
\midrule
White   & $10\,366\,482t$ & $7\,206\,360t$\\
Black   & $7\,460\,514t$  & $4\,893\,246t$\\
Yellow  & $4\,149\,387t$  & $5\,439\,213t$\\
Dappled & $7\,358\,060t$  & $3\,515\,820t$\\
\bottomrule
\end{tabular}
\end{center}
Substitution verifies all seven equations, and the eight coefficients
have gcd $1$. Therefore all eight counts are integers exactly when
$t$ is an integer: a B\'ezout combination of those counts equals $t$.
Every positive solution of the linear part has $t\ge1$, and the total
number of cattle is
\begin{equation}\label{eq:cattle-total-scale}
 T=50\,389\,082t.
\end{equation}

\subsection*{The square and the triangle produce Pell}
The square condition is
\[
 W+B=17\,826\,996t
 =4\cdot4\,456\,749t,
 \qquad4\,456\,749=3\cdot11\cdot29\cdot4657.
\]
The four factors in the last product are distinct primes (trial division
checks $4657$). Unique factorization therefore says that this is a square
exactly when
\[
 t=4\,456\,749r^2\qquad(r\ge1).
\]
Indeed, every prime in the displayed squarefree product must have odd
exponent in $t$, and every other prime must have even exponent.
Now the triangular condition becomes
\[
 11\,507\,447t=\frac{h(h+1)}2,
 \qquad
 51\,285\,802\,909\,803r^2=\frac{h(h+1)}2.
\]
Multiplying by $8$ and completing the square gives
\begin{equation}\label{eq:cattle-large-pell}
 \boxed{\quad X^2-410\,286\,423\,278\,424r^2=1,
 \qquad X=2h+1.\quad}
\end{equation}
Conversely, a positive solution of this equation has odd $X$, since the
coefficient of $r^2$ is divisible by $8$. Thus $h=(X-1)/2$ is an integer
and reconstructs the triangle. Together with the table, this proves an
equivalence, not just a necessary condition.

The coefficient is large, but contains a large square factor:
\[
 410\,286\,423\,278\,424=4\,729\,494\cdot9314^2.
\]
Put $Z=9314r$. The equivalent smaller-coefficient problem is
\begin{equation}\label{eq:cattle-small-pell}
 \boxed{\quad X^2-4\,729\,494Z^2=1,
 \qquad 9314\divides Z.\quad}
\end{equation}
Keeping the divisibility condition is essential: an arbitrary solution
of the smaller Pell equation need not give an integer $r$.

\subsection*{A short period, followed by a divisibility test}
The integer-state algorithm gives $a_0=2174$ and least period length
$L=92$ for $\sqrt{4\,729\,494}$. The final digit is $4348$; the preceding
$91$ digits form a palindrome. Since $92$ is even, the convergent
$p_{91}/q_{91}$ gives the fundamental solution
\begin{align*}
 u&=109931986732829734979866232821433543901088049,\\
 v&=50549485234315033074477819735540408986340.
\end{align*}
These values can be checked by the exact equality
$u^2-4\,729\,494v^2=1$. Write
$\varepsilon=u+v\sqrt{4\,729\,494}$ and
\[
 \varepsilon^j=X_j+Z_j\sqrt{4\,729\,494}.
\]
Proposition~\ref{prop:pell-powers} proves that this lists every positive
solution of the smaller Pell equation. We need the first $j$ for which
$9314\divides Z_j$.

We can find it without ever forming the enormous powers. Reduce the
multiplication recurrence modulo $9314$. Here
\[
 u\congr9063,\quad v\congr7708,\quad
 4\,729\,494\congr7296\pmod{9314}.
\]
The following finite computation checks every exponent up to the first
success, using only small integers:
\begin{quote}\begin{minipage}{\linewidth}\small
\begin{verbatim}
modulus = 9314
u, v, d = 9063, 7708, 7296
a, b = 1, 0
for j in range(1, 2330):
    a, b = ((u*a + d*v*b) % modulus,
            (v*a + u*b) % modulus)
    if b == 0:
        print(j, a, b)
        break
\end{verbatim}
\end{minipage}\end{quote}
Its output is \texttt{2329 9313 0}; no earlier exponent passes the test.
Thus $\varepsilon^{2329}$ produces the smallest admissible $r$ and
hence the smallest herd. This is a reproducible finite calculation,
distinct from the general Pell proofs above. Vardi explains a shorter
divisibility test using a quadratic finite field.

\subsection*{An enormous answer with a compact description}
Let $X+Z\sqrt{4\,729\,494}=\varepsilon^{2329}$. Then the least total is
\begin{equation}\label{eq:cattle-exact-total}
 T=50\,389\,082\cdot4\,456\,749
       \left(\frac{Z}{9314}\right)^2.
\end{equation}
This is an exact specification using integer arithmetic. The first
notebook computes $X,Z$ by repeated squaring, verifies the seven
proportions and both shape conditions, and checks
\[
 10^{206544}\le T<10^{206545}.
\]
The decimal expansion of the smallest herd's size therefore has
$206\,545$ digits:
\[
 T\approx7.760271406\times10^{206544}.
\]
Its size is a mathematical consequence of the conditions, rather than
a plausible census of animals.

Conjugation also gives a compact real-number expression, of the kind
developed in Vardi's exposition. Since
\[
 Z=\frac{\varepsilon^{2329}-\varepsilon^{-2329}}
          {2\sqrt{4\,729\,494}},
\]
substitution into \eqref{eq:cattle-exact-total} yields
\[
 T=\frac{25\,194\,541}{184\,119\,152}
       \left(\varepsilon^{4658}+\varepsilon^{-4658}-2\right).
\]
An expression can be short even when its decimal expansion is long.

\begin{bigidea}
The linear conditions leave one scale factor. Prime-exponent parity
makes that factor a fixed number times a square. The triangular
condition becomes Pell, and a divisibility test selects the allowed
powers of its fundamental solution. The same elementary tools that
solve $p^2-2q^2=1$ organize the cattle problem.
\end{bigidea}
